\chapter{Conclusion}
\label{sec:conclusion}

We have studied the problem of scheduling a chain of multi-mode tasks by an
end-to-end deadline in the presence of fail-stop failures. Even a
restricted version of this problem is NP-hard. In the one-PE case, 
dynamic programming makes it possible to precompute optimal scheduling
decisions for a time-discretized version of the problem under a
first-output execution policy. In the multi-PE case, we have proposed
several heuristics.  Using these heuristics in conjunction with temporal
redundancy and reactive rescheduling schemes yields a set of scheduling
strategies.  Simulation results over a large set of application
configurations show which scheduling strategy design choices are sound,
yielding a small set of top-performing strategies with different relative
effectiveness at different platform scales.


Our work can be extended in several ways.  Our temporal redundancy approach
could be extended to re-execute the application from a task $T_i$, $i > 1$,
and/or could be integrated as part of the dynamic program (at the cost of a
much larger state space).  Our experimental results provide a basis upon
which to develop more effective heuristics in the future, e.g., ones that
combine our dynamic programming approach with our greedy heuristics.
Another future work direction is to consider in-situ executions in which
consecutive tasks are merged (to avoid I/O).  Previously proposed task
merging strategies for the purpose of makespan minimization~\cite{6264675}
may be adapted to expected error minimization for multi-mode tasks. In the
multi-PE case, when tasks are parallel programs, an interesting problem is
to find a compromise between having each task use all PEs (as in
Section~\ref{sec:data_parallel}) and each task use a single PE (as in
Chapter~\ref{sec:multihost}).  Finally, a broader future direction is to go
beyond chains and consider more general task graphs.
